\section{Einleitung}
\label{sec:Einleitung}
Die folgende Projektdokumentation beschreibt den Ablauf des IHK-Abschlussprojektes, das im Rahmen der Ausbildung zum Fachinformatiker Fachrichtung Anwendungsentwicklung durchgeführt wurde. Der Ausbildungsbetrieb ist die MIBS AG, ein SAP-Beratungsunternehmen mit Sitz in Mülheim an der Ruhr. Zurzeit beschäftigt die MIBS AG rund 94 Mitarbeiter\footnote{Stand November 2024}. Das Unternehmen bietet SAP-Beratung und Projektmanagement an und vertreibt außerdem sein eigenes Produkt „SMART“, ein Testwerkzeug für SAP-Anwendungen.
Darüber hinaus entwickelt die MIBS AG ihr eigenes HCM-Tool\footnote{Software die Unternehmen bei der Verwaltung, Planung und Optimierung ihrer Personalressourcen unterstützt} namens \textbf{MIBS virtualDesk}. Dieses wird für die Organisation interner Prozesse, wie beispielsweise Urlaubsanträge, Krankmeldungen oder Zeitnachweise, genutzt.

\subsection{Projektbeschreibung}
\label{sec:Projektbeschreibung}
Derzeit werden im Backoffice der MIBS AG verschiedene operative Aufgaben manuell bearbeitet und in Excel-Tabellen dokumentiert. Hierzu gehören unter anderem das Übertragen von Abwesenheiten von Mitarbeitern aus dem MIBS virtualDesk sowie die Freigabe von Urlaubsanträgen nach deren Erfassung in der Excel-Datei.  
Dieses Dokument wird außerdem regelmäßig zur Planung interner Projekte und Schulungen verwendet. Auch für das Abgleichen von Verfügbarkeiten der Mitarbeiter durch Vorstand oder Teamleiter kommt dieses Dokument zum Einsatz.
In diesem Projekt soll sowohl der gerade beschriebene Prozess vereinfacht werden als auch die manuelle Pflege der Excel-Datei entfallen.


\subsection{Projektziel}
\label{sec:Projektziel}
Das Ziel des Projektes ist das Ersetzen der Excel-Datei durch eine Webanwendung innerhalb des MIBS virtualDesk der MIBS AG.
Folglich soll hierdurch der paralelle Pflegeaufwand entfallen.
Die Angestelltendaten sollen über das bestehende SAP-System verarbeitet werden.
Mitarbeiterabwesenheiten sollen aggregiert und über eine App innerhalb des MIBS virtualDesk verfügbar gemacht werden. Die Abwesenheitszeiträume sollen in einem Kalender, in dem jeder Mitarbeiter aufgelistet ist, ausgegeben werden. Des Weiteren sollen Funktionen für die Filterung der Mitarbeiter bereitgestellt werden.
Unter den angezeigten Daten sollen die Urlaubs- und Krankheitstage, sowie deren aktueller Status angezeigt werden. Zusätzlich sollen die Schulblöcke der Azubis ausgegeben werden.

\subsection{Projektumfeld}
\label{sec:Projektumfeld}
Auftraggeber des Projektes ist der Vorstand der MIBS AG, vertreten durch Marcel Ebert. Betreut wird das Projekt von der Abteilungsleitung des „MIBS Technology Hub“. Der zuständige Ansprechpartner ist Rafael Kopytto. Das Projekt ist für den Zeitraum 15.10.2024 – 21.11.2024 terminiert.

Die Abteilung MTH beobachtet und bewertet neue SAP Entwicklungstrends und unterstützt deren Kunden in der Implementierung dieser. Weiterhin betreut sie die Entwicklung und Wartung des MIBS virtualDesk.


Hierbei wurde sich auf die Entwicklung von RESTful-Anwendungen sowie darauf basierenden Webanwendungen spezialisiert. Primär kommen dabei ABAP RAP\footnote{\label{RAP}RESTful Application Programming – Modell von SAP zur Entwicklung cloudfähiger Anwendungen} für das Backend sowie SAPUI5\footnote{\label{fn:UI5}JavaScript-Framework von SAP zur Entwicklung responsiver Webanwendungen}, insbesondere das Framework Fiori Elements, für das Frontend zum Einsatz. 


\subsection{Projektbegründung}
\label{sec:Projektbegründung}
Die MIBS AG versucht fortlaufend ihre internen Prozesse zu optimieren. Besonders im Bereich des Backoffice wird versucht, manuelle und redundante Aufgaben zu automatisieren und umzustrukturieren. Der gesamte Prozess der Abwesenheitspflege erfordert einen kontinuierlichen manuellen Aufwand und beinhaltet ebenfalls eine redundante Datenpflege. Zum einen werden die Abwesenheiten der Mitarbeiterdaten täglich händisch in der Excel-Tabelle aktualisiert, welches einen wiederkehrenden Aufwand für das Backoffice bedeutet. Zum anderen besteht eine hohe Fehleranfälligkeit bei dem Übertragen der Daten.
Diese Umstände führten zu der Entscheidung, die Abwesenheitsübersicht als App im MIBS virtualDesk zu integrieren.

\subsection{Projektabgrenzung}
\label{sec:Projektabgrenzung}
Das Projekt ist so strukturiert, dass es ausschließlich als Datenanzeige fungiert. Die Hauptfunktion besteht darin, Informationen bereitzustellen, die aus einer definierten Datenquelle abgerufen werden. Es ist ausdrücklich festgelegt, dass die Anwendung keine Funktionen zur Änderung, Löschung oder Neuanlage von Daten bereitstellt. Die zugrunde liegenden Daten bleiben unverändert.

